%%%%%%%%%%%%%%%%%%%%%%%%% CHAPTER - 7 %%%%%%%%%%%%%%%%%%%%%%%%%
\chapter{Summary and Conclusions}
\label{C7}
%%%%%%%%%%%%%%%%%%%%%%%%%%%%%%%%%%%%%%%%%%%%%%%%%%%%%%%%%%%%%%%
\par This thesis addresses some important issues in three-flavor neutrino oscillation in the context of upcoming high-precision long-baseline experiments. Neutrinos are omnipresent and they are the second most abundant particles in the Universe after photons. There are a variety of neutrino sources (both natural and artificial) producing neutrinos over a wide range of energies in the range of $10^{-4}$ to $10^{18}$ eV or so and traveling distances from a few meters to hundreds of Gpc. Neutrinos were postulated first by Wolfgang Pauli in 1930 to explain how beta-decay could conserve energy, momentum, and angular momentum. In 1934, Enrico Fermi named the new particle as ``neutrino", which means the ``little neutral one" in the Italian language. In 1956, Clyde L. Cowan and Frederick Reines first observed neutrinos in a nuclear reactor. There are three (3) light active flavor neutrinos: $\nu_e$, $\nu_\mu$, and $\nu_\tau$ - experimentally confirmed by the precision data on the \textit{Z}-decay width at the $e^+e^-$ collider at LEP \cite{ALEPH:2005ab}. They only take part in weak interactions through massive \textit{W}$^\pm$ and \textit{Z}$^0$ gauge bosons. They are electrically neutral, left-handed, and form \textit{SU}$(2)$ doublets in the Standard Model with the corresponding left-handed charged leptons. In the basic SM of particle physics, neutrinos are massless due to the absence of right-handed neutrinos. \\
\par Over the past two decades, data from several world-class neutrino experiments firmly established that neutrinos change flavor after propagating a finite distance in space and time - a phenomenon known as ``neutrino oscillation". Neutrino oscillation demands non-zero neutrino mass and mixing. Hence, non-zero neutrino mass is the first experimental proof for physics beyond the Standard Model for which the Nobel Prize was awarded to Takaaki Kajita and Arthur McDonald in 2015. Neutrino oscillation, in the $3\nu$ paradigm, is mainly governed by six oscillation parameters, viz., solar mixing angle ($\theta_{12}$), atmospheric mixing angle ($\theta_{23}$), reactor mixing angle ($\theta_{13}$), Dirac CP-violating phase ($\delta_{\mathrm{CP}}$), solar mass-squared splitting ($\Delta m^2_{21} \equiv m^2_{2}-m^2_{1}$), and atmospheric mass-squared splitting ($\Delta m^2_{31} \equiv m^2_{3}-m^2_{1}$). Still, there are a few fundamental issues that need to be resolved in neutrino oscillation experiments, namely, i) the neutrino mass ordering (whether $\Delta m^2_{31}>0$, known as normal mass ordering - NMO, or $\Delta m^2_{31}<0$, termed as inverted mass ordering - IMO), ii) the octant of $\theta_{23}$ (whether $\theta_{23}<45^\circ$ known as lower octant - LO, or $\theta_{23}>45^\circ$ termed as higher octant - HO), and iii) whether the CP is violated in the neutrino sector like quark sector and if so, then what is the precise value of the CP phase? The next-generation high-precision long-baseline (LBL) experiments are capable to address these issues at a high confidence level. \\\vspace{0.1cm}
In this thesis, we study i) the physics reach of the upcoming LBL experiment, Deep Underground Neutrino Experiment (DUNE) in the US, to establish the possible deviation from maximal mixing of $\theta_{23}$, ii) to achieve an improved precision on the 2-3 oscillation parameters using the synergy between the DUNE and Tokai-to-Hyper-Kamiokande (T2HK) experiments, and iii) the possible impact of the flavor-dependent long-range neutrino interactions on the measurements of neutrino oscillation parameters in DUNE and T2HK.\\
\section{Establishing the deviation from maximal mixing of $\theta_{23}$ in DUNE:}\hspace{0.2cm}The robustness of the three-flavor neutrino oscillation paradigm got established on a strong footing after the discovery of non-zero $\theta_{13}$ in the Daya Bay experiment and opened the door for searching CP violation in currently running and upcoming LBL experiments. The recent global fit studies ~\cite{Capozzi:2021fjo,deSalas:2020pgw,Esteban:2020cvm,NuFIT} of the three-flavor neutrino oscillation parameters indicate that $\sin^2\theta_{23}\neq0.5$ or $\sin^22\theta_{23}\neq1$. This gives rise to the issue of non-maximal $\theta_{23}$, which is quite important from the neutrino mass-mixing models point of view. In this work, we study the prospect of the next-generation long-baseline (LBL) experiment DUNE to establish the non-maximal $\theta_{23}$ in light of the current neutrino oscillation data.\\
DUNE \cite{DUNE:2020jqi} is an upcoming LBL experiment with a baseline of 1285 km where neutrinos travel from Fermilab to Homestake Mine in South Dakota. It will use an on-axis, wide-band neutrino beam with high intensity covering both first and second oscillation maxima. Some of its notable features are the excellent energy resolution of the Liquid Argon TPC (LArTPC) having 40 kt fiducial volume, large matter effect due to a longer baseline, and less systematic uncertainties in the appearance channel. The neutrino flux peaks around the first oscillation maxima which occurs at 2.5 GeV for 1285 km.\\ 
The study of non-maximal $\theta_{23}$ is mostly driven by the disappearance channel, and hence the uncertainty in $\delta_{\mathrm{CP}}$ does not affect our results much. As the disappearance channel is statistics-driven, it also helps to improve the precision on $\theta_{23}$. Our benchmark choices of the oscillation parameters are from Capozzi $et$ al. \cite{Capozzi:2021fjo}, where we have only considered the NMO as the true mass ordering in Nature. In our benchmark choices, $\sin^2\theta_{23} = 0.455$ (in LO), and $\delta_{\mathrm{CP}} = 223^\circ$ (belongs to lower-half-plane). We fix the value of $\theta_{12},\, \theta_{13}$, and $\Delta m^2_{21}$ as they are already measured with high precision using solar and reactor experiments. In our analysis, we have allowed the values of $\theta_{23}, \Delta m^2_{31}$, and $\delta_{\mathrm{CP}}$ to vary within their $3\sigma$ allowed range as mentioned in \cite{Capozzi:2021fjo}.\\
As the disappearance probability $\propto$ $\sin^22\theta_{23}$, degeneracies persist at the total event level. However, due to the superior energy resolution of DUNE, this degeneracy can be lifted using spectral information alongwith the total event rates. Some energy bins around the 1st oscillation maxima ($i.e.$, at 2.5 GeV) show the maximum sensitivity towards the non-maximal $\theta_{23}$ whereas, the performance of lower and higher energy bins deteriorates for the benchmark value of $\Delta m^2_{31}$. We show, with [3.5 yr $\nu$ + 3.5 yr $\bar{\nu}$] run, DUNE can establish non-maximal $\theta_{23}$ at $\sim$ $5\sigma$ C.L. with total rate + shape analyses whereas, it gets dropped to $1.5\sigma$ C.L. with an analysis based on only total event rates.\\
\par We see that with 7 years of run time equally shared by neutrino and antineutrino, DUNE can establish non-maximal $\theta_{23}$ with $\sim$ 4.2$\sigma$ C.L. at the benchmark choice of true $\sin^2\theta_{23}$. It is obvious that the more the true (benchmark) value of $\sin^2\theta_{23}$ shifts towards the maximal mixing value, the more it will be difficult to establish non-maximality of $\theta_{23}$. It is noteworthy to mention that if the true value of $\sin^2\theta_{23}$ shifts towards the upper range of the $1\sigma$ uncertainty of $\sin^2\theta_{23}$ ($i.e.$, maximum limit within $1\sigma$ tolerance in the vicinity of MM), DUNE can still establish non-maximality of $\sin^2\theta_{23}$ with $\sim 2.1\sigma$ C.L. However, this study is improved for the enhanced run time [$i.e.$, 5 yr $\nu$ + 5 yr $\bar{\nu}$] as it is mainly governed by the statistics. We observe that a $3\sigma\, (5\sigma)$ determination of non-maximal $\theta_{23}$ is possible in DUNE with an exposure of 336 kt$\cdot$MW$\cdot$years if true $\sin^2\theta_{23}\lesssim 0.465\, (0.450)$ or $\sin^2\theta_{23}\gtrsim0.554\, (0.572)$ for any value of true $\delta_{\mathrm{CP}}$ in its present $3\sigma$ range and true NMO.\\
\section{Improved precision on the 2-3 oscillation parameters using the synergy between DUNE and T2HK:} \hspace{0.2cm} A high-precision measurement of $\Delta m^2_{31}$ and $\theta_{23}$ is inevitable to estimate the Earth's matter effect in long-baseline experiments which in turn plays an important role in addressing the issue of neutrino mass ordering and to measure the value of CP phase in $3\nu$ framework. Following the footprints of present-generation measurements and near-future sensitivities from different oscillation experiments, our work examines how next-generation experiments, DUNE and T2HK, bring improvement in the 2-3 sector which will also help to decipher the long-standing flavor problem glorifying the oscillation-based neutrino research in the precision era. We highlight the relevance of the complementarity between DUNE and T2HK in determining the sensitivity towards deviation from maximal $\theta_{23}$, excluding the wrong octant of $\theta_{23}$, and establishing the precision on atmospheric parameters.\\
\par The next-generation long baseline experiment in Japan, T2HK, \cite{Hyper-KamiokandeProto-:2015xww, Hyper-Kamiokande:2018ofw} will use an off-axis, narrow-band neutrino beam that will travel through a distance of 295 km from J-PARC to Hyper-Kamiokande. This experiment has several features that are complementary to DUNE. Due to its short baseline, T2HK will experience less matter effect than DUNE and can have better prospects in measuring intrinsic CP violation due to $\delta_{\mathrm{CP}}$. T2HK's larger detector volume than DUNE can improve any disappearance-driven phenomena with the power of more enriched statistics. To compensate for the low antineutrino cross section, T2HK has been planned with 3 times higher run time in antineutrino mode than neutrino [2.5 yr $\nu$ + 7.5 yr $\bar{\nu}$]. T2HK's larger systematic uncertainty in appearance channel ($5\%$) than DUNE ($2\%$) is responsible for reduced sensitivity in any analysis driven by appearance channel, whereas, its improved disappearance systematics ($3.5\%$) than DUNE ($5\%$) along with larger detector volume helps to improve any disappearance-driven sensitivity.\\
We show that T2HK, with its benchmark exposure of 2431 kt$\cdot$MW$\cdot$yr, can establish non-maximal $\theta_{23}$ at $5.6\sigma$ C.L. which deteriorates a bit for DUNE to $5\sigma$ with its benchmark exposure of 480 kt$\cdot$MW$\cdot$yr at the benchmark choice of $\sin^2\theta_{23}$ ($i.e.$, 0.455) \cite{Capozzi:2021fjo} assuming NMO. T2HK's enriched statistics and less systematic uncertainty in the disappearance channel enhance its performance over DUNE. The performance is much improved ($\sim 7.6\sigma$) in the combined setup of DUNE + T2HK as compared to their standalone setups. The notable finding is that, in Nature, if true $\sin^2\theta_{23}$ attains the upper boundary of the current $1\sigma$ uncertainty (0.473) which is close to MM, only DUNE + T2HK can achieve $3\sigma$ sensitivity of non-maximal $\theta_{23}$ with the present best-fit values of oscillation parameters. Based on the current benchmark choices of oscillation parameters \cite{Capozzi:2021fjo}, the range of true values of $\sin^{2}\theta_{23}$ that can be differentiated from MM choices, by DUNE + T2HK with just half of their nominal exposures, cannot be achieved by either of the individual experiments even with their respective full projected exposures.\\
We study one of the most pressing unsolved issues in $3\nu$ paradigm, $i.e.$, the octant ambiguity of $\theta_{23}$. As it is driven by appearance channel, DUNE performs better ($5.2\sigma$) in excluding the wrong octant solution of $\theta_{23}$ with its nominal exposure than T2HK ($4.4\sigma$) at benchmark value of $\sin^2\theta_{23} = 0.455$ for true $\delta_{\mathrm{CP}} = 223^\circ$ and true NMO. However, the joint setup of DUNE + T2HK outperforms ($7.4\sigma$) in our study than their solo performances. Here also, our finding show that the joint performance of DUNE + T2HK can give $\sim 4\sigma$ sensitivity to our study at the upper edge of $1\sigma$ uncertainty in $\sin^2\theta_{23}$ ($i.e.$, in the vicinity of MM). At lower confidence, T2HK wins due to larger statistics, whereas, at higher confidence, DUNE wins due to lesser systematics in the appearance channel. One more noteworthy finding in this context is that, with just 0.25 times the benchmark exposure of the individual experiments, the combined setup can exclude $60\%$ of $\theta_{23}$ from the wrong octant solutions at $3\sigma$ C.L.\\
\par We address one more pertinent issue in the precision domain of neutrino oscillation, viz., precision measurements of $\sin^2\theta_{23}$ and $\Delta m^2_{31}$ around their benchmark choices ($i.e.$, 0.455 and $2.522\times10^{-3}$ eV$^2$) for NMO. Due to the statistical abundance, the standalone T2HK shows more precise measurements than DUNE and, of course, their combined setup enhances the performance in these measurements than their individual setups at their nominal exposures. The highlighted result from this section is that, the combined setup improves the present achievable precision of $\sin^2\theta_{23}$ and $\Delta m^2_{31}$, by a factor of 7 and 5, respectively. While standalone DUNE and T2HK with low exposures ($\sim 0.25$ times nominal exposure) cannot rule out clone solutions in $\sin^2\theta_{23}$ at $3\sigma$, the combined DUNE + T2HK provide degeneracy-free measurements. In the case of $\Delta m^2_{31}$, using approximately $20\%$ of the individual exposures of DUNE and T2HK together can achieve an impressive relative $1\sigma$ precision of $0.25\%$.\\
The combination of DUNE and T2HK can exclude the HO only in antineutrino
mode at $3\sigma$ C.L. breaking $(\sin^2\theta_{23} - \delta_{\mathrm{CP}})$ degeneracy due to higher $\bar{\nu}$ statistics in T2HK. So, the majority of the appearance events are free from fake (matter-induced) CP-phase. However, for the combined neutrino and antineutrino modes, we can give more precise measurements of $\sin^2\theta_{23}$ and $\Delta m^2_{31}$ by squeezing the parameter space irrespective of all experimental setups under discussion. DUNE's excellent energy-resolution helps to precisely measure $\Delta m^2_{31}$ whereas, T2HK's statistical enrichment assists to precisely measure $\sin^2\theta_{23}$. We also show the reduction in the parameter space $(\sin^2\theta_{23} - \delta_{\mathrm{CP}})$ using DUNE and T2HK, either in isolation or in combination. We find that, at half-exposure, the synergy between DUNE and T2HK is the only hope to exclude the MM solution and disfavor the CP-conserving region.\\
\section{Impact of long-range neutrino interactions on the measurements of oscillation parameters at DUNE and T2HK :} \hspace{0.2cm} Neutrino oscillation, having the bright prospect of measuring oscillation parameters with high precision in the next-generation experiments, may unravel novel signatures of BSM physics. The symmetries of the Standard Model are preserved under the gauge group \textit{SU}$(3)_C\,\otimes\,$\textit{SU}$(2)_L\,\otimes\,$\textit{U}$(1)_Y$. Now, a new, anomaly-free interaction can be added through a simple extension of the gauge group, viz., \textit{U}$(1)^\prime$, which brings some new symmetries in the Standard Model. These symmetries are chosen as the difference between the lepton numbers, $i.e.$, $(L_e-L_\mu)$, $(L_e-L_\tau)$, and $(L_\mu-L_\tau)$, manifesting the new interactions in Nature. Now, if the mass of the corresponding gauge boson $Z^\prime$ is very light ($\leq\,10^{-10}$ eV), the new interaction is long-ranged and characterized as a long-range interaction (LRI). The next-generation LBL experiments DUNE and Hyper-K may feel the presence of these new interactions \cite{Singh:2023nek}. In this work, we showcase the prominent effect of LRI along with the standard neutrino interaction (SI) in DUNE, T2HK, and their combined setups on measuring non-maximal $\theta_{23}$ and leptonic CP violation. The value of $(\Delta m^2_{31}/2E)$ for DUNE at 1st oscillation maxima is $\sim$ $10^{-13}$ eV which is compatible with the order of magnitude of LRI potential and hence DUNE and T2HK are good portals in probing the existence of LRI potential through oscillation. We assume NMO as a true mass ordering in Nature and the true LRI potential as $1.4\times10^{-14}$ eV, $1.0\times10^{-14}$ eV, and $0.73\times10^{-14}$ eV for $(L_e-L_\mu)$, $(L_e-L_\tau)$, and $(L_\mu-L_\tau)$ symmetries, respectively \cite{Singh:2023nek} throughout our analyses. In general, $L_e-L_\mu$ and $L_e-L_\tau$ are sensitive to $\delta_{\mathrm{CP}}$ and $\theta_{23}$ whereas, 
$L_\mu-L_\tau$ symmetry is  more sensitive to $\Delta m^2_{31}$ and $\theta_{23}$.\\
When the strength of long-range interactions increases ($V_{\rm{LRI}} > 10^{-14}$ eV), the distinguishable fraction of $\sin^{2}\theta_{23}$ values in Nature, deviating from maximal mixing, increases significantly. This enhancement is more pronounced under the $L_{\mu} - L_{\tau}$ symmetry, where the combined capabilities of DUNE and T2HK can achieve a coverage of approximately 90\% if $V_{\mu\tau}\,(\mathrm{true}\sim5\times10^{-14})$ eV. This improvement is attributed to the substantial disappearance event statistics in both DUNE and T2HK. For an illustrative true LRI value of $\sim 10^{-14}$ eV, we find that if the true octant of $\sin^{2}\theta_{23}$ in Nature corresponds to the lower octant (LO), the presence of LRI interactions under all the three symmetries reduces the sensitivity to deviations from maximal mixing compared to the SI scenario. Conversely, the sensitivity improves when $\sin^{2}\theta_{23}$ resides in the higher octant (HO) and is close to maximal mixing. We observe similar findings when studying the sensitivity towards the exclusion of the wrong octant. $L_{\mu} - L_{\tau}$ symmetry affects the most. We also investigate the impact of subdominant LRI on establishing leptonic CP violation (CPV). Our analysis reveals that if subdominant LRI exists in Nature, the fraction of $\delta_{\rm{CP}}$ values capable of establishing CPV at $\geq 3\sigma$ under each symmetry decreases significantly. This reduction arises from the critical role of $\sin^{2}2\theta_{13}$ evolution as a function of neutrino energy in the presence of matter. For an illustrative value of $V_{\rm{LRI}} \sim 10^{-14}$ eV, we find that the presence of LRI has a relatively greater impact on CPV sensitivity in DUNE compared to T2HK. This can be attributed to the higher prevalence of fake CP asymmetry in DUNE due to its longer baseline. Additionally, the $L_\mu - L_\tau$ symmetry introduces the most significant degradation in CPV sensitivity. This deterioration is primarily driven by the degeneracy between $V_{\rm{LRI}}$ and the uncertainty in $\sin^{2}\theta_{23}$. We also analyze the impact of LRI on the allowed parameter ranges in $(\sin^2\theta_{23}-\delta_{\mathrm{CP}})$ plane and observe varying effects across different symmetries. Under the $L_{e}-L_{\tau}$ symmetry, the allowed ranges remain largely consistent with those under SI, as sensitivity in this case is predominantly driven by the appearance statistics. For the $L_{e}-L_{\mu}$ symmetry, a noticeable deterioration in the precision of $\sin^{2}\theta_{23}$ is observed in DUNE when the data is generated with $\sin^{2}\theta_{23}$ in the LO. However, the complementarity between DUNE and T2HK resolves these degeneracies, resulting in stringent allowed regions under LRI that are comparable to those in the SI scenario. To conclude, our findings emphasize that while subdominant LRIs may influence the sensitivity and precision measurements of oscillation parameters, the complementary capabilities of DUNE and T2HK mitigate these challenges to a large extent. This synergy underscores the importance of a combined experimental approach in disentangling the impact of LRIs from neutrino oscillation under SI.
